% ============================================================
% PHBS 硕士学位论文配置文件
% ============================================================
% public 模板的唯一 metadata 源是根 configs.tex；
% parts/*/configs.tex 只作为同步出来的镜像，不单独维护。
%
% 阶段说明 (修改下方 \stage 的值):
%   blind   - 盲审版: 隐藏学生姓名、学号、导师信息，无致谢
%   defense - 答辩版: 显示学生信息，隐藏导师信息，无致谢
%   final   - 最终版: 显示所有信息，包含致谢
%
% 也可以通过 make 命令指定阶段:
%   make blind    - 编译盲审版
%   make defense  - 编译答辩版
%   make final    - 编译最终版
%   make          - 使用下方配置的阶段
% ============================================================

% ============================================================
% 【重要】当前阶段设置
% ============================================================
% 可选值: blind / defense / final
\newcommand{\stage}{final}
% 构建系统会在工作目录里自动生成 stage.auto.tex 覆盖阶段值；
% 用户手工直接编译某个分册时，则回退到上面的默认值。
\IfFileExists{stage.auto.tex}{\input{stage.auto.tex}}{}

% ============================================================
% 【重要】版式与参考文献模式
% ============================================================
% 字体配置:
%   strict_pku - 交稿模式，优先使用 Windows 宋体系字体
%   portable   - 兼容模式，允许后备字体链
\newcommand{\fontprofile}{strict_pku}

% 参考文献体例:
%   author_year - 著者-出版年制（默认）
%   numeric     - 顺序编码制
% 具体 biblatex style 由 parts/en|zh/main.tex 绑定为 gb7714-2015/gb7714-2015ay；
% PHBS 附件 1 不强制 bibliography style，若未来有明确要求，可只切换该绑定。
\newcommand{\bibmode}{author_year}

% ============================================================
% 论文标题
% ============================================================
% 中文标题 (使用 \zhquote{} 输入中文引号)
\newcommand{\zhtitle}{\zhquote{摸鱼}的经济学原理——信息不对称下的最优偷懒策略}
% 英文标题 (使用 \titleenquote{} 输入标题中的英文引号)
\newcommand{\entitle}{The Economics of \titleenquote{Idleness} - Optimal Laziness under Asymmetric Information}
% 标题行数 (用于封面排版，根据标题长度调整)
\newcommand{\titlelines}{3}

% ============================================================
% 作者信息
% ============================================================
\newcommand{\zhauthorname}{李嘉图}
\newcommand{\enauthorname}{Li Jiatu}
\newcommand{\mystudentid}{2201212000}

% ============================================================
% 导师信息
% ============================================================
% 格式: 姓名 + 两个空格 + 职称
% 职称: 教授 / 副教授 / 讲师 / 助理教授
\newcommand{\zhmentorname}{亚当·斯密\ \ 教授}
% 英文格式: Prof. / A.P. (Associate Professor) / Lec. (Lecturer)
\newcommand{\enmentorname}{Prof.\ Adam Smith}

% ============================================================
% 专业信息
% ============================================================
% 专业名称
% 可选: 金融硕士 / 西方经济学 / 企业管理 / 新闻与传播硕士
\newcommand{\zhmajor}{金融硕士}
% 可选: Master of Finance / Master of Economics / Master of Management / Master of Journalism and Communication
\newcommand{\enmajor}{Master of Finance}
% 研究方向 (金融专业填写，其他专业留空)
% 可选: 金融管理 / 数量金融 / 金融科技 / 金融投资 / 财经传媒
\newcommand{\majordirection}{金融管理}
% 学位类型: true = 学术学位, false = 专业学位
\newcommand{\isacademicdegree}{false}

% ============================================================
% 关键词
% ============================================================
% 中文关键词可用中文或英文逗号分隔，模板会统一转成中文逗号
\newcommand{\zhkeywords}{信息不对称, 偷懒, 激励, 博弈论}
% 英文关键词可用中文或英文逗号分隔，模板会统一转成半角逗号并让每个关键词首字母大写
\newcommand{\enkeywords}{Asymmetric information, Laziness, Incentive, Game theory}

% ============================================================
% 日期
% ============================================================
% 具体时间以教务为准: 初稿3月，送审4月，答辩5月，最终6月
\newcommand{\theyear}{2025}
\newcommand{\themonth}{5}

% ============================================================
% 以下内容自动处理，无需修改
% ============================================================

% 阶段判断逻辑
\usepackage{ifthen}

% 定义阶段布尔变量
\newif\ifblind
\newif\ifdefense
\newif\iffinal

% 根据 \stage 设置布尔变量
\ifthenelse{\equal{\stage}{blind}}{
	\blindtrue \defensefalse \finalfalse
}{
	\ifthenelse{\equal{\stage}{defense}}{
		\blindfalse \defensetrue \finalfalse
	}{
		\blindfalse \defensefalse \finaltrue
	}
}

% 是否显示学生信息 (defense 和 final 显示)
\newif\ifshowstudent
\ifblind \showstudentfalse \else \showstudenttrue \fi

% 是否显示导师信息 (仅 final 显示)
\newif\ifshowmentor
\iffinal \showmentortrue \else \showmentorfalse \fi

% 是否显示致谢 (仅 final 显示)
\newif\ifshowacknowledgement
\iffinal \showacknowledgementtrue \else \showacknowledgementfalse \fi

% 根据阶段设置实际显示的作者/导师信息
\ifshowstudent
	\newcommand{\zhauthor}{\zhauthorname}
	\newcommand{\enauthor}{\enauthorname}
	\newcommand{\thestudentid}{\mystudentid}
\else
	\newcommand{\zhauthor}{***}
	\newcommand{\enauthor}{***}
	\newcommand{\thestudentid}{**********}
\fi

\ifshowmentor
	\newcommand{\zhmentor}{\zhmentorname}
	\newcommand{\enmentor}{\enmentorname}
\else
	\newcommand{\zhmentor}{******}
	\newcommand{\enmentor}{******}
\fi

% ============================================================
% 额外的包和自定义命令
% ============================================================
\usepackage{tikz}
\usetikzlibrary{arrows.meta, positioning, calc}
